\subsection{Topologies on ordered vector spaces}

Let E be an ordered vector space. A topology on E is compatible with the ordered vector space structure of E if it is both compatible with the vector space structure of E and subject to the following axiom :

(TO) The convex cone of the \(x\) with \(x \geqslant 0\) , is closed in E.

An ordered vector space E with a compatible topology is called an ordered topological vector space.

Examples. --- The space \(R^{n}\) with its usual topology and the order structure that is the product of the order structure of its factors is an ordered topological vector space. On the other hand, for \(n \geqslant 2\) , when \(R^{n}\) carries the lexicographical order (S, III, § 2.6), the usual topology is not compatible with the ordered vector space structure of \(R^{n}\) .

Let \(\mathbf{A}\) be a set; the vector space \(\mathcal{B}(\mathbf{A};\mathbf{R})\) of real valued bounded functions defined on \(\mathbf{A}\) , with the topology defined by the norm \(\| x\| = \sup_{t\in \mathbf{A}}|x(t)|\) and the order structure

induced by the product order structure of \(R^{A}\) , is an ordered topological vector space.

In an ordered topological vector space E, the set of elements \(x \leqslant 0\) is closed; since translations are homeomorphisms, we deduce that, for all \(a \in E\) , the set of elements \(x \geqslant a\) (resp. \(x \leqslant a\) ) is closed. Since \(\{0\}\) is the intersection of the sets \(x \geqslant 0\) and \(x \leqslant 0\) , it follows that \(\{0\}\) is closed and that E is Hausdorff.

PROPOSITION 18. --- In an ordered topological vector space E, let H be a set directed by the relation \(\leqslant\) . If the section filter of H has a limit in E, then this limit is the upper bound of H.

For, let \(b = \lim_{x\in \mathbf{H}}x\) ; for every \(y\in \mathbf{H}\) , the set of \(x\in \mathbf{H}\) such that \(x\geqslant y\) is a set of

the section filter of H, therefore b is a cluster point of this set; but as the set \(x \geqslant y\) is closed in E, we have \(b \geqslant y\) , thus b is an upper bound of H. On the other hand, if a is an upper bound of H, then H is contained in the closed set \(x \leqslant a\) ; as b is a cluster point of H, we have \(b \leqslant a\) , which completes the proof (II, p.~72, exerc. 42).

